% Problem 1 -- included from hw.tex via % Problem 1 -- included from hw.tex via % Problem 1 -- included from hw.tex via % Problem 1 -- included from hw.tex via \input{problems/problem1}.
% Paths (figures/, *.py) are relative to the directory of hw.tex.

%===============================================================================
\section*{Problem 1: Reachability Analysis}
%===============================================================================

\noindent\textbf{Conventions.}
For the autonomous system, $\trj{\cdot}$ denotes the unique trajectory that starts from
$\state_0$ at time $\tvar_0$. For the controlled system, $\trju{\ctrl}{\cdot}$ denotes the
trajectory generated by the control signal $\ctrl(\cdot)\in\Usig$, the set of admissible control
signals on $[\tvar_0,\tvar]$, which we assume to be nonempty. Trajectories are continuous in
time. Since ``$A\subseteq B$'' means ``$\state_0\in A \Rightarrow \state_0\in B$'', each proof
below takes an arbitrary point of the left-hand set and exhibits the witnesses required by the
definition of the right-hand set, and each counterexample exhibits a point that lies in the
left-hand set but not in the right-hand set.

\subsection*{Part 1: autonomous system $\dot\state=\dyn(\state)$}

\begin{enumerate}[label=(\alph*), itemsep=3pt]

\item \TRUE\
Let $\state_0\in\BRS{\tvar}$, i.e.\ $\trj{\tvar}\in\target$. Choose $\taux:=\tvar$. Then
$\taux\in[\tvar_0,\tvar]$ and $\trj{\taux}\in\target$, which is exactly the defining condition of
$\BRT{\tvar}{\target}$; hence $\state_0\in\BRT{\tvar}{\target}$. In words, the tube collects the
states that visit $\target$ at \emph{some} time in the horizon, in particular those that are in
$\target$ exactly at the final time.

\item \TRUE\
Let $\state_0\in\BRT{\tvar}{\target}$: there is $\taux\in[\tvar_0,\tvar]$ with
$\trj{\taux}\in\target$. Because $\tvar<\tvar'$ we have $[\tvar_0,\tvar]\subseteq[\tvar_0,\tvar']$,
so the same $\taux$ satisfies $\taux\in[\tvar_0,\tvar']$ and $\trj{\taux}\in\target$; therefore
$\state_0\in\BRT{\tvar'}{\target}$. Backward-reachable tubes are monotone (non-decreasing) in the
horizon.

\item \FALSE\
The backward-reachable \emph{set} only constrains the state at time $\tvar$; it does not care
whether the trajectory passed through $\failure$ on the way. Counterexample ($n=1$):
\[
\dot\state = 1,\qquad \tvar_0=0,\quad \tvar=2,\qquad \target=\{2\},\quad \failure=\{1\}
\]
(both nonempty and disjoint). From $\state_0=0$ the trajectory is $\trj{\taux}=\taux$, so
$\trj{2}=2\in\target$ and $\state_0\in\BRS{2}$. However, the only time in $[0,2]$ at which the
trajectory is in $\target$ is $\taux=2$, and for $s=1\in[\tvar_0,\taux]$ we have
$\trj{1}=1\in\failure$. The avoid condition fails for the only admissible $\taux$, so
$\state_0\notin\RAT{2}$. Hence $\BRS{\tvar}\not\subseteq\RAT{\tvar}$ in general.

\item \TRUE\
Let $\state_0\in\RAT{\tvar}$. By definition there is $\taux\in[\tvar_0,\tvar]$ such that
$\trj{\taux}\in\target$ \emph{and} $\trj{s}\notin\failure$ for all $s\in[\tvar_0,\taux]$. The
first conjunct, with the same witness $\taux$, is precisely the defining condition of
$\BRT{\tvar}{\target}$; hence $\state_0\in\BRT{\tvar}{\target}$. Dropping a conjunct from a
membership condition can only enlarge the set, so the reach-avoid tube is always contained in the
reach tube.

\end{enumerate}

\subsection*{Part 2: controlled system $\dot\state=\dyn(\state,\ctrl)$}

\begin{enumerate}[label=(\alph*), itemsep=3pt]

\item \TRUE\
Let $\state_0\in\iBRT{\tvar}{\target}$: for \emph{every} $\ctrl(\cdot)\in\Usig$ there is
$\taux_{\ctrl}\in[\tvar_0,\tvar]$ with $\trju{\ctrl}{\taux_{\ctrl}}\in\target$. Since
$\Usig\neq\emptyset$ we may fix any particular $\bar\ctrl(\cdot)\in\Usig$; then
$\taux_{\bar\ctrl}\in[\tvar_0,\tvar]$ and $\trju{\bar\ctrl}{\taux_{\bar\ctrl}}\in\target$, so the
pair $(\bar\ctrl,\taux_{\bar\ctrl})$ witnesses $\state_0\in\eBRT{\tvar}{\target}$. In short,
``for all $\ctrl$'' implies ``there exists $\ctrl$'' because the set of admissible controls is
nonempty: if reaching $\target$ cannot be avoided, then it can certainly be enforced.

\item \FALSE\
The enforceable reach-avoid tube is
\[
\eRAT{\tvar} := \bigl\{\state_0\in\reals^n :\ \exists\,\ctrl(\cdot)\in\Usig,\ \exists\,\taux\in[\tvar_0,\tvar],\
\trju{\ctrl}{\taux}\in\target \ \wedge\ \forall s\in[\tvar_0,\taux],\ \trju{\ctrl}{s}\notin\failure\bigr\}.
\]
The intersection $\eBRT{\tvar}{\target}\cap\Safe{\tvar}{\failure}$ allows \emph{one} control to
reach $\target$ and a \emph{different} control to avoid $\failure$, whereas $\eRAT{\tvar}$ requires
a \emph{single} control to do both. Counterexample ($n=1$):
\[
\dot\state=\ctrl,\quad \ctrl\in[-1,1],\qquad \tvar_0=0,\quad \tvar=2,\qquad
\target=\{2\},\quad \failure=\{1\},\qquad \state_0=0.
\]
\begin{itemize}[itemsep=0pt]
\item $\state_0\in\eBRT{2}{\target}$: with $\ctrl\equiv1$, $\trju{\ctrl}{\taux}=\taux$, so
      $\trju{\ctrl}{2}=2\in\target$.
\item $\state_0\in\Safe{2}{\failure}$: with $\ctrl\equiv0$, $\trju{\ctrl}{\taux}\equiv0\notin\failure$
      for all $\taux\in[0,2]$.
\item $\state_0\notin\eRAT{2}$: if some control satisfies $\trju{\ctrl}{\taux}=2$ for some $\taux$,
      the (continuous) trajectory goes from $0$ to $2$, so by the intermediate value theorem
      $\trju{\ctrl}{s}=1\in\failure$ for some $s\in(0,\taux)$. The avoid condition therefore fails
      for every candidate pair $(\ctrl,\taux)$.
\end{itemize}
Hence $\eBRT{\tvar}{\target}\cap\Safe{\tvar}{\failure}\not\subseteq\eRAT{\tvar}$, and the equality
is false. (The reverse inclusion fails in general as well, because $\eRAT{\tvar}$ only asks to avoid
$\failure$ \emph{until} $\target$ is reached, while $\Safe{\tvar}{\failure}$ asks to avoid it over
the whole horizon. With $\dot\state=1+\ctrl$, $\ctrl\in[-\tfrac12,\tfrac12]$, $\tvar_0=0$, $\tvar=10$,
$\target=\{1\}$, $\failure=\{2\}$, the state $\state_0=0$ lies in $\eRAT{10}$, since $\ctrl\equiv0$
reaches $1$ at $\taux=1$ without touching $2$, but not in $\Safe{10}{\failure}$, since
$\dot\state\ge\tfrac12$ forces $\trju{\ctrl}{4}\ge2$ and thus every trajectory crosses $2$ before
time $4<10$.)

\item \FALSE\
Only the inclusion ``$\subseteq$'' holds in general.

\emph{Proof of $\Safe{\tvar}{\failure_1\cup\failure_2}\subseteq\Safe{\tvar}{\failure_1}\cap\Safe{\tvar}{\failure_2}$.}
If $\state_0\in\Safe{\tvar}{\failure_1\cup\failure_2}$ with witness $\ctrl(\cdot)$, then
$\trju{\ctrl}{\taux}\notin\failure_1\cup\failure_2$ for all $\taux\in[\tvar_0,\tvar]$, i.e.\
$\trju{\ctrl}{\taux}\notin\failure_1$ and $\trju{\ctrl}{\taux}\notin\failure_2$. The same
$\ctrl(\cdot)$ therefore witnesses both $\state_0\in\Safe{\tvar}{\failure_1}$ and
$\state_0\in\Safe{\tvar}{\failure_2}$.

\emph{Counterexample to ``$\supseteq$''.} The controls that avoid $\failure_1$ and $\failure_2$ may be
different, and no single control may avoid both. Take $n=2$, $\state=(\state_1,\state_2)$,
\[
\dot\state_1=1,\quad \dot\state_2=\ctrl,\quad \ctrl\in[-1,1],\qquad \tvar_0=0,\quad \tvar=1,\qquad
\state_0=(0,0),
\]
\[
\failure_1=\{(1,\state_2):\state_2\ge0\},\qquad
\failure_2=\{(1,\state_2):\state_2\le0\},\qquad
\failure_1\cup\failure_2=\{(1,\state_2):\state_2\in\reals\}.
\]
\begin{itemize}[itemsep=0pt]
\item $\state_0\in\Safe{1}{\failure_1}$: with $\ctrl\equiv-1$ the trajectory is $(\taux,-\taux)$. It has
      $\state_1=1$ only at $\taux=1$, where $\state_2=-1<0$, so it never enters $\failure_1$.
\item $\state_0\in\Safe{1}{\failure_2}$: symmetrically, $\ctrl\equiv+1$ gives $(\taux,\taux)$, which
      never enters $\failure_2$.
\item $\state_0\notin\Safe{1}{\failure_1\cup\failure_2}$: $\state_1(\taux)=\taux$ regardless of the
      control, so at $\taux=1\in[0,1]$ every trajectory satisfies
      $\trju{\ctrl}{1}=(1,\state_2(1))\in\failure_1\cup\failure_2$.
\end{itemize}
Thus $\Safe{\tvar}{\failure_1}\cap\Safe{\tvar}{\failure_2}\supsetneq\Safe{\tvar}{\failure_1\cup\failure_2}$
in general: being able to dodge each hazard separately does not imply being able to dodge both at
once. Taking complements (note $\reals^n\setminus\Safe{\tvar}{\failure}=\iBRT{\tvar}{\failure}$),
the inevitable BRT of a union of failure sets can be \emph{strictly larger} than the union of the
individual inevitable BRTs; Problem~3 exhibits exactly this effect.

\end{enumerate}
.
% Paths (figures/, *.py) are relative to the directory of hw.tex.

%===============================================================================
\section*{Problem 1: Reachability Analysis}
%===============================================================================

\noindent\textbf{Conventions.}
For the autonomous system, $\trj{\cdot}$ denotes the unique trajectory that starts from
$\state_0$ at time $\tvar_0$. For the controlled system, $\trju{\ctrl}{\cdot}$ denotes the
trajectory generated by the control signal $\ctrl(\cdot)\in\Usig$, the set of admissible control
signals on $[\tvar_0,\tvar]$, which we assume to be nonempty. Trajectories are continuous in
time. Since ``$A\subseteq B$'' means ``$\state_0\in A \Rightarrow \state_0\in B$'', each proof
below takes an arbitrary point of the left-hand set and exhibits the witnesses required by the
definition of the right-hand set, and each counterexample exhibits a point that lies in the
left-hand set but not in the right-hand set.

\subsection*{Part 1: autonomous system $\dot\state=\dyn(\state)$}

\begin{enumerate}[label=(\alph*), itemsep=3pt]

\item \TRUE\
Let $\state_0\in\BRS{\tvar}$, i.e.\ $\trj{\tvar}\in\target$. Choose $\taux:=\tvar$. Then
$\taux\in[\tvar_0,\tvar]$ and $\trj{\taux}\in\target$, which is exactly the defining condition of
$\BRT{\tvar}{\target}$; hence $\state_0\in\BRT{\tvar}{\target}$. In words, the tube collects the
states that visit $\target$ at \emph{some} time in the horizon, in particular those that are in
$\target$ exactly at the final time.

\item \TRUE\
Let $\state_0\in\BRT{\tvar}{\target}$: there is $\taux\in[\tvar_0,\tvar]$ with
$\trj{\taux}\in\target$. Because $\tvar<\tvar'$ we have $[\tvar_0,\tvar]\subseteq[\tvar_0,\tvar']$,
so the same $\taux$ satisfies $\taux\in[\tvar_0,\tvar']$ and $\trj{\taux}\in\target$; therefore
$\state_0\in\BRT{\tvar'}{\target}$. Backward-reachable tubes are monotone (non-decreasing) in the
horizon.

\item \FALSE\
The backward-reachable \emph{set} only constrains the state at time $\tvar$; it does not care
whether the trajectory passed through $\failure$ on the way. Counterexample ($n=1$):
\[
\dot\state = 1,\qquad \tvar_0=0,\quad \tvar=2,\qquad \target=\{2\},\quad \failure=\{1\}
\]
(both nonempty and disjoint). From $\state_0=0$ the trajectory is $\trj{\taux}=\taux$, so
$\trj{2}=2\in\target$ and $\state_0\in\BRS{2}$. However, the only time in $[0,2]$ at which the
trajectory is in $\target$ is $\taux=2$, and for $s=1\in[\tvar_0,\taux]$ we have
$\trj{1}=1\in\failure$. The avoid condition fails for the only admissible $\taux$, so
$\state_0\notin\RAT{2}$. Hence $\BRS{\tvar}\not\subseteq\RAT{\tvar}$ in general.

\item \TRUE\
Let $\state_0\in\RAT{\tvar}$. By definition there is $\taux\in[\tvar_0,\tvar]$ such that
$\trj{\taux}\in\target$ \emph{and} $\trj{s}\notin\failure$ for all $s\in[\tvar_0,\taux]$. The
first conjunct, with the same witness $\taux$, is precisely the defining condition of
$\BRT{\tvar}{\target}$; hence $\state_0\in\BRT{\tvar}{\target}$. Dropping a conjunct from a
membership condition can only enlarge the set, so the reach-avoid tube is always contained in the
reach tube.

\end{enumerate}

\subsection*{Part 2: controlled system $\dot\state=\dyn(\state,\ctrl)$}

\begin{enumerate}[label=(\alph*), itemsep=3pt]

\item \TRUE\
Let $\state_0\in\iBRT{\tvar}{\target}$: for \emph{every} $\ctrl(\cdot)\in\Usig$ there is
$\taux_{\ctrl}\in[\tvar_0,\tvar]$ with $\trju{\ctrl}{\taux_{\ctrl}}\in\target$. Since
$\Usig\neq\emptyset$ we may fix any particular $\bar\ctrl(\cdot)\in\Usig$; then
$\taux_{\bar\ctrl}\in[\tvar_0,\tvar]$ and $\trju{\bar\ctrl}{\taux_{\bar\ctrl}}\in\target$, so the
pair $(\bar\ctrl,\taux_{\bar\ctrl})$ witnesses $\state_0\in\eBRT{\tvar}{\target}$. In short,
``for all $\ctrl$'' implies ``there exists $\ctrl$'' because the set of admissible controls is
nonempty: if reaching $\target$ cannot be avoided, then it can certainly be enforced.

\item \FALSE\
The enforceable reach-avoid tube is
\[
\eRAT{\tvar} := \bigl\{\state_0\in\reals^n :\ \exists\,\ctrl(\cdot)\in\Usig,\ \exists\,\taux\in[\tvar_0,\tvar],\
\trju{\ctrl}{\taux}\in\target \ \wedge\ \forall s\in[\tvar_0,\taux],\ \trju{\ctrl}{s}\notin\failure\bigr\}.
\]
The intersection $\eBRT{\tvar}{\target}\cap\Safe{\tvar}{\failure}$ allows \emph{one} control to
reach $\target$ and a \emph{different} control to avoid $\failure$, whereas $\eRAT{\tvar}$ requires
a \emph{single} control to do both. Counterexample ($n=1$):
\[
\dot\state=\ctrl,\quad \ctrl\in[-1,1],\qquad \tvar_0=0,\quad \tvar=2,\qquad
\target=\{2\},\quad \failure=\{1\},\qquad \state_0=0.
\]
\begin{itemize}[itemsep=0pt]
\item $\state_0\in\eBRT{2}{\target}$: with $\ctrl\equiv1$, $\trju{\ctrl}{\taux}=\taux$, so
      $\trju{\ctrl}{2}=2\in\target$.
\item $\state_0\in\Safe{2}{\failure}$: with $\ctrl\equiv0$, $\trju{\ctrl}{\taux}\equiv0\notin\failure$
      for all $\taux\in[0,2]$.
\item $\state_0\notin\eRAT{2}$: if some control satisfies $\trju{\ctrl}{\taux}=2$ for some $\taux$,
      the (continuous) trajectory goes from $0$ to $2$, so by the intermediate value theorem
      $\trju{\ctrl}{s}=1\in\failure$ for some $s\in(0,\taux)$. The avoid condition therefore fails
      for every candidate pair $(\ctrl,\taux)$.
\end{itemize}
Hence $\eBRT{\tvar}{\target}\cap\Safe{\tvar}{\failure}\not\subseteq\eRAT{\tvar}$, and the equality
is false. (The reverse inclusion fails in general as well, because $\eRAT{\tvar}$ only asks to avoid
$\failure$ \emph{until} $\target$ is reached, while $\Safe{\tvar}{\failure}$ asks to avoid it over
the whole horizon. With $\dot\state=1+\ctrl$, $\ctrl\in[-\tfrac12,\tfrac12]$, $\tvar_0=0$, $\tvar=10$,
$\target=\{1\}$, $\failure=\{2\}$, the state $\state_0=0$ lies in $\eRAT{10}$, since $\ctrl\equiv0$
reaches $1$ at $\taux=1$ without touching $2$, but not in $\Safe{10}{\failure}$, since
$\dot\state\ge\tfrac12$ forces $\trju{\ctrl}{4}\ge2$ and thus every trajectory crosses $2$ before
time $4<10$.)

\item \FALSE\
Only the inclusion ``$\subseteq$'' holds in general.

\emph{Proof of $\Safe{\tvar}{\failure_1\cup\failure_2}\subseteq\Safe{\tvar}{\failure_1}\cap\Safe{\tvar}{\failure_2}$.}
If $\state_0\in\Safe{\tvar}{\failure_1\cup\failure_2}$ with witness $\ctrl(\cdot)$, then
$\trju{\ctrl}{\taux}\notin\failure_1\cup\failure_2$ for all $\taux\in[\tvar_0,\tvar]$, i.e.\
$\trju{\ctrl}{\taux}\notin\failure_1$ and $\trju{\ctrl}{\taux}\notin\failure_2$. The same
$\ctrl(\cdot)$ therefore witnesses both $\state_0\in\Safe{\tvar}{\failure_1}$ and
$\state_0\in\Safe{\tvar}{\failure_2}$.

\emph{Counterexample to ``$\supseteq$''.} The controls that avoid $\failure_1$ and $\failure_2$ may be
different, and no single control may avoid both. Take $n=2$, $\state=(\state_1,\state_2)$,
\[
\dot\state_1=1,\quad \dot\state_2=\ctrl,\quad \ctrl\in[-1,1],\qquad \tvar_0=0,\quad \tvar=1,\qquad
\state_0=(0,0),
\]
\[
\failure_1=\{(1,\state_2):\state_2\ge0\},\qquad
\failure_2=\{(1,\state_2):\state_2\le0\},\qquad
\failure_1\cup\failure_2=\{(1,\state_2):\state_2\in\reals\}.
\]
\begin{itemize}[itemsep=0pt]
\item $\state_0\in\Safe{1}{\failure_1}$: with $\ctrl\equiv-1$ the trajectory is $(\taux,-\taux)$. It has
      $\state_1=1$ only at $\taux=1$, where $\state_2=-1<0$, so it never enters $\failure_1$.
\item $\state_0\in\Safe{1}{\failure_2}$: symmetrically, $\ctrl\equiv+1$ gives $(\taux,\taux)$, which
      never enters $\failure_2$.
\item $\state_0\notin\Safe{1}{\failure_1\cup\failure_2}$: $\state_1(\taux)=\taux$ regardless of the
      control, so at $\taux=1\in[0,1]$ every trajectory satisfies
      $\trju{\ctrl}{1}=(1,\state_2(1))\in\failure_1\cup\failure_2$.
\end{itemize}
Thus $\Safe{\tvar}{\failure_1}\cap\Safe{\tvar}{\failure_2}\supsetneq\Safe{\tvar}{\failure_1\cup\failure_2}$
in general: being able to dodge each hazard separately does not imply being able to dodge both at
once. Taking complements (note $\reals^n\setminus\Safe{\tvar}{\failure}=\iBRT{\tvar}{\failure}$),
the inevitable BRT of a union of failure sets can be \emph{strictly larger} than the union of the
individual inevitable BRTs; Problem~3 exhibits exactly this effect.

\end{enumerate}
.
% Paths (figures/, *.py) are relative to the directory of hw.tex.

%===============================================================================
\section*{Problem 1: Reachability Analysis}
%===============================================================================

\noindent\textbf{Conventions.}
For the autonomous system, $\trj{\cdot}$ denotes the unique trajectory that starts from
$\state_0$ at time $\tvar_0$. For the controlled system, $\trju{\ctrl}{\cdot}$ denotes the
trajectory generated by the control signal $\ctrl(\cdot)\in\Usig$, the set of admissible control
signals on $[\tvar_0,\tvar]$, which we assume to be nonempty. Trajectories are continuous in
time. Since ``$A\subseteq B$'' means ``$\state_0\in A \Rightarrow \state_0\in B$'', each proof
below takes an arbitrary point of the left-hand set and exhibits the witnesses required by the
definition of the right-hand set, and each counterexample exhibits a point that lies in the
left-hand set but not in the right-hand set.

\subsection*{Part 1: autonomous system $\dot\state=\dyn(\state)$}

\begin{enumerate}[label=(\alph*), itemsep=3pt]

\item \TRUE\
Let $\state_0\in\BRS{\tvar}$, i.e.\ $\trj{\tvar}\in\target$. Choose $\taux:=\tvar$. Then
$\taux\in[\tvar_0,\tvar]$ and $\trj{\taux}\in\target$, which is exactly the defining condition of
$\BRT{\tvar}{\target}$; hence $\state_0\in\BRT{\tvar}{\target}$. In words, the tube collects the
states that visit $\target$ at \emph{some} time in the horizon, in particular those that are in
$\target$ exactly at the final time.

\item \TRUE\
Let $\state_0\in\BRT{\tvar}{\target}$: there is $\taux\in[\tvar_0,\tvar]$ with
$\trj{\taux}\in\target$. Because $\tvar<\tvar'$ we have $[\tvar_0,\tvar]\subseteq[\tvar_0,\tvar']$,
so the same $\taux$ satisfies $\taux\in[\tvar_0,\tvar']$ and $\trj{\taux}\in\target$; therefore
$\state_0\in\BRT{\tvar'}{\target}$. Backward-reachable tubes are monotone (non-decreasing) in the
horizon.

\item \FALSE\
The backward-reachable \emph{set} only constrains the state at time $\tvar$; it does not care
whether the trajectory passed through $\failure$ on the way. Counterexample ($n=1$):
\[
\dot\state = 1,\qquad \tvar_0=0,\quad \tvar=2,\qquad \target=\{2\},\quad \failure=\{1\}
\]
(both nonempty and disjoint). From $\state_0=0$ the trajectory is $\trj{\taux}=\taux$, so
$\trj{2}=2\in\target$ and $\state_0\in\BRS{2}$. However, the only time in $[0,2]$ at which the
trajectory is in $\target$ is $\taux=2$, and for $s=1\in[\tvar_0,\taux]$ we have
$\trj{1}=1\in\failure$. The avoid condition fails for the only admissible $\taux$, so
$\state_0\notin\RAT{2}$. Hence $\BRS{\tvar}\not\subseteq\RAT{\tvar}$ in general.

\item \TRUE\
Let $\state_0\in\RAT{\tvar}$. By definition there is $\taux\in[\tvar_0,\tvar]$ such that
$\trj{\taux}\in\target$ \emph{and} $\trj{s}\notin\failure$ for all $s\in[\tvar_0,\taux]$. The
first conjunct, with the same witness $\taux$, is precisely the defining condition of
$\BRT{\tvar}{\target}$; hence $\state_0\in\BRT{\tvar}{\target}$. Dropping a conjunct from a
membership condition can only enlarge the set, so the reach-avoid tube is always contained in the
reach tube.

\end{enumerate}

\subsection*{Part 2: controlled system $\dot\state=\dyn(\state,\ctrl)$}

\begin{enumerate}[label=(\alph*), itemsep=3pt]

\item \TRUE\
Let $\state_0\in\iBRT{\tvar}{\target}$: for \emph{every} $\ctrl(\cdot)\in\Usig$ there is
$\taux_{\ctrl}\in[\tvar_0,\tvar]$ with $\trju{\ctrl}{\taux_{\ctrl}}\in\target$. Since
$\Usig\neq\emptyset$ we may fix any particular $\bar\ctrl(\cdot)\in\Usig$; then
$\taux_{\bar\ctrl}\in[\tvar_0,\tvar]$ and $\trju{\bar\ctrl}{\taux_{\bar\ctrl}}\in\target$, so the
pair $(\bar\ctrl,\taux_{\bar\ctrl})$ witnesses $\state_0\in\eBRT{\tvar}{\target}$. In short,
``for all $\ctrl$'' implies ``there exists $\ctrl$'' because the set of admissible controls is
nonempty: if reaching $\target$ cannot be avoided, then it can certainly be enforced.

\item \FALSE\
The enforceable reach-avoid tube is
\[
\eRAT{\tvar} := \bigl\{\state_0\in\reals^n :\ \exists\,\ctrl(\cdot)\in\Usig,\ \exists\,\taux\in[\tvar_0,\tvar],\
\trju{\ctrl}{\taux}\in\target \ \wedge\ \forall s\in[\tvar_0,\taux],\ \trju{\ctrl}{s}\notin\failure\bigr\}.
\]
The intersection $\eBRT{\tvar}{\target}\cap\Safe{\tvar}{\failure}$ allows \emph{one} control to
reach $\target$ and a \emph{different} control to avoid $\failure$, whereas $\eRAT{\tvar}$ requires
a \emph{single} control to do both. Counterexample ($n=1$):
\[
\dot\state=\ctrl,\quad \ctrl\in[-1,1],\qquad \tvar_0=0,\quad \tvar=2,\qquad
\target=\{2\},\quad \failure=\{1\},\qquad \state_0=0.
\]
\begin{itemize}[itemsep=0pt]
\item $\state_0\in\eBRT{2}{\target}$: with $\ctrl\equiv1$, $\trju{\ctrl}{\taux}=\taux$, so
      $\trju{\ctrl}{2}=2\in\target$.
\item $\state_0\in\Safe{2}{\failure}$: with $\ctrl\equiv0$, $\trju{\ctrl}{\taux}\equiv0\notin\failure$
      for all $\taux\in[0,2]$.
\item $\state_0\notin\eRAT{2}$: if some control satisfies $\trju{\ctrl}{\taux}=2$ for some $\taux$,
      the (continuous) trajectory goes from $0$ to $2$, so by the intermediate value theorem
      $\trju{\ctrl}{s}=1\in\failure$ for some $s\in(0,\taux)$. The avoid condition therefore fails
      for every candidate pair $(\ctrl,\taux)$.
\end{itemize}
Hence $\eBRT{\tvar}{\target}\cap\Safe{\tvar}{\failure}\not\subseteq\eRAT{\tvar}$, and the equality
is false. (The reverse inclusion fails in general as well, because $\eRAT{\tvar}$ only asks to avoid
$\failure$ \emph{until} $\target$ is reached, while $\Safe{\tvar}{\failure}$ asks to avoid it over
the whole horizon. With $\dot\state=1+\ctrl$, $\ctrl\in[-\tfrac12,\tfrac12]$, $\tvar_0=0$, $\tvar=10$,
$\target=\{1\}$, $\failure=\{2\}$, the state $\state_0=0$ lies in $\eRAT{10}$, since $\ctrl\equiv0$
reaches $1$ at $\taux=1$ without touching $2$, but not in $\Safe{10}{\failure}$, since
$\dot\state\ge\tfrac12$ forces $\trju{\ctrl}{4}\ge2$ and thus every trajectory crosses $2$ before
time $4<10$.)

\item \FALSE\
Only the inclusion ``$\subseteq$'' holds in general.

\emph{Proof of $\Safe{\tvar}{\failure_1\cup\failure_2}\subseteq\Safe{\tvar}{\failure_1}\cap\Safe{\tvar}{\failure_2}$.}
If $\state_0\in\Safe{\tvar}{\failure_1\cup\failure_2}$ with witness $\ctrl(\cdot)$, then
$\trju{\ctrl}{\taux}\notin\failure_1\cup\failure_2$ for all $\taux\in[\tvar_0,\tvar]$, i.e.\
$\trju{\ctrl}{\taux}\notin\failure_1$ and $\trju{\ctrl}{\taux}\notin\failure_2$. The same
$\ctrl(\cdot)$ therefore witnesses both $\state_0\in\Safe{\tvar}{\failure_1}$ and
$\state_0\in\Safe{\tvar}{\failure_2}$.

\emph{Counterexample to ``$\supseteq$''.} The controls that avoid $\failure_1$ and $\failure_2$ may be
different, and no single control may avoid both. Take $n=2$, $\state=(\state_1,\state_2)$,
\[
\dot\state_1=1,\quad \dot\state_2=\ctrl,\quad \ctrl\in[-1,1],\qquad \tvar_0=0,\quad \tvar=1,\qquad
\state_0=(0,0),
\]
\[
\failure_1=\{(1,\state_2):\state_2\ge0\},\qquad
\failure_2=\{(1,\state_2):\state_2\le0\},\qquad
\failure_1\cup\failure_2=\{(1,\state_2):\state_2\in\reals\}.
\]
\begin{itemize}[itemsep=0pt]
\item $\state_0\in\Safe{1}{\failure_1}$: with $\ctrl\equiv-1$ the trajectory is $(\taux,-\taux)$. It has
      $\state_1=1$ only at $\taux=1$, where $\state_2=-1<0$, so it never enters $\failure_1$.
\item $\state_0\in\Safe{1}{\failure_2}$: symmetrically, $\ctrl\equiv+1$ gives $(\taux,\taux)$, which
      never enters $\failure_2$.
\item $\state_0\notin\Safe{1}{\failure_1\cup\failure_2}$: $\state_1(\taux)=\taux$ regardless of the
      control, so at $\taux=1\in[0,1]$ every trajectory satisfies
      $\trju{\ctrl}{1}=(1,\state_2(1))\in\failure_1\cup\failure_2$.
\end{itemize}
Thus $\Safe{\tvar}{\failure_1}\cap\Safe{\tvar}{\failure_2}\supsetneq\Safe{\tvar}{\failure_1\cup\failure_2}$
in general: being able to dodge each hazard separately does not imply being able to dodge both at
once. Taking complements (note $\reals^n\setminus\Safe{\tvar}{\failure}=\iBRT{\tvar}{\failure}$),
the inevitable BRT of a union of failure sets can be \emph{strictly larger} than the union of the
individual inevitable BRTs; Problem~3 exhibits exactly this effect.

\end{enumerate}
.
% Paths (figures/, *.py) are relative to the directory of hw.tex.

%===============================================================================
\section*{Problem 1: Reachability Analysis}
%===============================================================================

\noindent\textbf{Conventions.}
For the autonomous system, $\trj{\cdot}$ denotes the unique trajectory that starts from
$\state_0$ at time $\tvar_0$. For the controlled system, $\trju{\ctrl}{\cdot}$ denotes the
trajectory generated by the control signal $\ctrl(\cdot)\in\Usig$, the set of admissible control
signals on $[\tvar_0,\tvar]$, which we assume to be nonempty. Trajectories are continuous in
time. Since ``$A\subseteq B$'' means ``$\state_0\in A \Rightarrow \state_0\in B$'', each proof
below takes an arbitrary point of the left-hand set and exhibits the witnesses required by the
definition of the right-hand set, and each counterexample exhibits a point that lies in the
left-hand set but not in the right-hand set.

\subsection*{Part 1: autonomous system $\dot\state=\dyn(\state)$}

\begin{enumerate}[label=(\alph*), itemsep=3pt]

\item \TRUE\
Let $\state_0\in\BRS{\tvar}$, i.e.\ $\trj{\tvar}\in\target$. Choose $\taux:=\tvar$. Then
$\taux\in[\tvar_0,\tvar]$ and $\trj{\taux}\in\target$, which is exactly the defining condition of
$\BRT{\tvar}{\target}$; hence $\state_0\in\BRT{\tvar}{\target}$. In words, the tube collects the
states that visit $\target$ at \emph{some} time in the horizon, in particular those that are in
$\target$ exactly at the final time.

\item \TRUE\
Let $\state_0\in\BRT{\tvar}{\target}$: there is $\taux\in[\tvar_0,\tvar]$ with
$\trj{\taux}\in\target$. Because $\tvar<\tvar'$ we have $[\tvar_0,\tvar]\subseteq[\tvar_0,\tvar']$,
so the same $\taux$ satisfies $\taux\in[\tvar_0,\tvar']$ and $\trj{\taux}\in\target$; therefore
$\state_0\in\BRT{\tvar'}{\target}$. Backward-reachable tubes are monotone (non-decreasing) in the
horizon.

\item \FALSE\
The backward-reachable \emph{set} only constrains the state at time $\tvar$; it does not care
whether the trajectory passed through $\failure$ on the way. Counterexample ($n=1$):
\[
\dot\state = 1,\qquad \tvar_0=0,\quad \tvar=2,\qquad \target=\{2\},\quad \failure=\{1\}
\]
(both nonempty and disjoint). From $\state_0=0$ the trajectory is $\trj{\taux}=\taux$, so
$\trj{2}=2\in\target$ and $\state_0\in\BRS{2}$. However, the only time in $[0,2]$ at which the
trajectory is in $\target$ is $\taux=2$, and for $s=1\in[\tvar_0,\taux]$ we have
$\trj{1}=1\in\failure$. The avoid condition fails for the only admissible $\taux$, so
$\state_0\notin\RAT{2}$. Hence $\BRS{\tvar}\not\subseteq\RAT{\tvar}$ in general.

\item \TRUE\
Let $\state_0\in\RAT{\tvar}$. By definition there is $\taux\in[\tvar_0,\tvar]$ such that
$\trj{\taux}\in\target$ \emph{and} $\trj{s}\notin\failure$ for all $s\in[\tvar_0,\taux]$. The
first conjunct, with the same witness $\taux$, is precisely the defining condition of
$\BRT{\tvar}{\target}$; hence $\state_0\in\BRT{\tvar}{\target}$. Dropping a conjunct from a
membership condition can only enlarge the set, so the reach-avoid tube is always contained in the
reach tube.

\end{enumerate}

\subsection*{Part 2: controlled system $\dot\state=\dyn(\state,\ctrl)$}

\begin{enumerate}[label=(\alph*), itemsep=3pt]

\item \TRUE\
Let $\state_0\in\iBRT{\tvar}{\target}$: for \emph{every} $\ctrl(\cdot)\in\Usig$ there is
$\taux_{\ctrl}\in[\tvar_0,\tvar]$ with $\trju{\ctrl}{\taux_{\ctrl}}\in\target$. Since
$\Usig\neq\emptyset$ we may fix any particular $\bar\ctrl(\cdot)\in\Usig$; then
$\taux_{\bar\ctrl}\in[\tvar_0,\tvar]$ and $\trju{\bar\ctrl}{\taux_{\bar\ctrl}}\in\target$, so the
pair $(\bar\ctrl,\taux_{\bar\ctrl})$ witnesses $\state_0\in\eBRT{\tvar}{\target}$. In short,
``for all $\ctrl$'' implies ``there exists $\ctrl$'' because the set of admissible controls is
nonempty: if reaching $\target$ cannot be avoided, then it can certainly be enforced.

\item \FALSE\
The enforceable reach-avoid tube is
\[
\eRAT{\tvar} := \bigl\{\state_0\in\reals^n :\ \exists\,\ctrl(\cdot)\in\Usig,\ \exists\,\taux\in[\tvar_0,\tvar],\
\trju{\ctrl}{\taux}\in\target \ \wedge\ \forall s\in[\tvar_0,\taux],\ \trju{\ctrl}{s}\notin\failure\bigr\}.
\]
The intersection $\eBRT{\tvar}{\target}\cap\Safe{\tvar}{\failure}$ allows \emph{one} control to
reach $\target$ and a \emph{different} control to avoid $\failure$, whereas $\eRAT{\tvar}$ requires
a \emph{single} control to do both. Counterexample ($n=1$):
\[
\dot\state=\ctrl,\quad \ctrl\in[-1,1],\qquad \tvar_0=0,\quad \tvar=2,\qquad
\target=\{2\},\quad \failure=\{1\},\qquad \state_0=0.
\]
\begin{itemize}[itemsep=0pt]
\item $\state_0\in\eBRT{2}{\target}$: with $\ctrl\equiv1$, $\trju{\ctrl}{\taux}=\taux$, so
      $\trju{\ctrl}{2}=2\in\target$.
\item $\state_0\in\Safe{2}{\failure}$: with $\ctrl\equiv0$, $\trju{\ctrl}{\taux}\equiv0\notin\failure$
      for all $\taux\in[0,2]$.
\item $\state_0\notin\eRAT{2}$: if some control satisfies $\trju{\ctrl}{\taux}=2$ for some $\taux$,
      the (continuous) trajectory goes from $0$ to $2$, so by the intermediate value theorem
      $\trju{\ctrl}{s}=1\in\failure$ for some $s\in(0,\taux)$. The avoid condition therefore fails
      for every candidate pair $(\ctrl,\taux)$.
\end{itemize}
Hence $\eBRT{\tvar}{\target}\cap\Safe{\tvar}{\failure}\not\subseteq\eRAT{\tvar}$, and the equality
is false. (The reverse inclusion fails in general as well, because $\eRAT{\tvar}$ only asks to avoid
$\failure$ \emph{until} $\target$ is reached, while $\Safe{\tvar}{\failure}$ asks to avoid it over
the whole horizon. With $\dot\state=1+\ctrl$, $\ctrl\in[-\tfrac12,\tfrac12]$, $\tvar_0=0$, $\tvar=10$,
$\target=\{1\}$, $\failure=\{2\}$, the state $\state_0=0$ lies in $\eRAT{10}$, since $\ctrl\equiv0$
reaches $1$ at $\taux=1$ without touching $2$, but not in $\Safe{10}{\failure}$, since
$\dot\state\ge\tfrac12$ forces $\trju{\ctrl}{4}\ge2$ and thus every trajectory crosses $2$ before
time $4<10$.)

\item \FALSE\
Only the inclusion ``$\subseteq$'' holds in general.

\emph{Proof of $\Safe{\tvar}{\failure_1\cup\failure_2}\subseteq\Safe{\tvar}{\failure_1}\cap\Safe{\tvar}{\failure_2}$.}
If $\state_0\in\Safe{\tvar}{\failure_1\cup\failure_2}$ with witness $\ctrl(\cdot)$, then
$\trju{\ctrl}{\taux}\notin\failure_1\cup\failure_2$ for all $\taux\in[\tvar_0,\tvar]$, i.e.\
$\trju{\ctrl}{\taux}\notin\failure_1$ and $\trju{\ctrl}{\taux}\notin\failure_2$. The same
$\ctrl(\cdot)$ therefore witnesses both $\state_0\in\Safe{\tvar}{\failure_1}$ and
$\state_0\in\Safe{\tvar}{\failure_2}$.

\emph{Counterexample to ``$\supseteq$''.} The controls that avoid $\failure_1$ and $\failure_2$ may be
different, and no single control may avoid both. Take $n=2$, $\state=(\state_1,\state_2)$,
\[
\dot\state_1=1,\quad \dot\state_2=\ctrl,\quad \ctrl\in[-1,1],\qquad \tvar_0=0,\quad \tvar=1,\qquad
\state_0=(0,0),
\]
\[
\failure_1=\{(1,\state_2):\state_2\ge0\},\qquad
\failure_2=\{(1,\state_2):\state_2\le0\},\qquad
\failure_1\cup\failure_2=\{(1,\state_2):\state_2\in\reals\}.
\]
\begin{itemize}[itemsep=0pt]
\item $\state_0\in\Safe{1}{\failure_1}$: with $\ctrl\equiv-1$ the trajectory is $(\taux,-\taux)$. It has
      $\state_1=1$ only at $\taux=1$, where $\state_2=-1<0$, so it never enters $\failure_1$.
\item $\state_0\in\Safe{1}{\failure_2}$: symmetrically, $\ctrl\equiv+1$ gives $(\taux,\taux)$, which
      never enters $\failure_2$.
\item $\state_0\notin\Safe{1}{\failure_1\cup\failure_2}$: $\state_1(\taux)=\taux$ regardless of the
      control, so at $\taux=1\in[0,1]$ every trajectory satisfies
      $\trju{\ctrl}{1}=(1,\state_2(1))\in\failure_1\cup\failure_2$.
\end{itemize}
Thus $\Safe{\tvar}{\failure_1}\cap\Safe{\tvar}{\failure_2}\supsetneq\Safe{\tvar}{\failure_1\cup\failure_2}$
in general: being able to dodge each hazard separately does not imply being able to dodge both at
once. Taking complements (note $\reals^n\setminus\Safe{\tvar}{\failure}=\iBRT{\tvar}{\failure}$),
the inevitable BRT of a union of failure sets can be \emph{strictly larger} than the union of the
individual inevitable BRTs; Problem~3 exhibits exactly this effect.

\end{enumerate}
